% GENERATED by scripts/manual_tables.py from the adapter registry. Do not edit:
% rerun the script. Hand-written text: blocks/preprocessing/wavelet_bands/
\ubblock{preprocessing.wavelet_bands}%
  {Wavelet bands}% what the app calls it
  {Wavelet bands (stationary wavelet transform)}% registry display name
  {preprocess}% insert-modal tab
  {Signal}% input
  {Signal}% output
  {Adapters/preprocessing\_wavelet\_bands.py}% code
  {preprocessing/wavelet_bands}% hand-written folder under blocks/
\ubblockdescription{preprocessing.wavelet_bands}{Stationary (undecimated) wavelet transform: splits the span into octave layers, each about half the speed of the one above, all sample-aligned with the recording and adding up to it. One layer --- the `level` you choose --- goes on down the chain as an ordinary signal; the block page draws every layer with its frequency range in Hz.}
\ubblockflags{preprocessing.wavelet_bands}{%
}
\ubblocksideinputs{preprocessing.wavelet_bands}{%
}
\ubblockparams{preprocessing.wavelet_bands}{%
  \ubparam{wavelet}{\texttt{db4}}{one of \texttt{db4}, \texttt{db2}, \texttt{db8}, \texttt{sym4}, \texttt{sym8}, \texttt{coif2}, \texttt{haar}}{The wavelet: db4 is a compact, moderately smooth default; haar is the sharpest, sym8/db8/coif2 the smoothest.}%
  \ubparam{levels}{0}{whole number from 0 to 20}{How many octave layers to split the span into. 0 = chosen from the span and the sample rate: deep enough to reach 0.001 Hz, as deep as the span allows.}%
  \ubparam{level}{4}{whole number from 0 to 20}{Which layer goes on down the chain: 1 = the fastest octave (fs/4 to fs/2), each next level half as fast; 0 = the residual approximation (everything slower than the deepest layer). Every layer is drawn on the block page with its Hz range.}%
}
